\clearpage
\fancyhead[L]{\small BRAMBILLA. EFFEMERIDI ASTRONOMICHE.}
\SourceLeaf{634}
\begin{center}\small Göttingische gelehrte Anzeigen. Stück 64. Seite 633..637. 1823 April 21.\\[3pt]\rule{.72\linewidth}{.3pt}\end{center}

\emph{Effemeridi Astronomiche di Milano per l'anno \textup{1823} calcolate da} \textsc{Enrico Brambilla}. \emph{Con Appendice.} Mailand 1822. Die Ephemeriden selbst 112 Seiten, der Anhang 88 Seiten Octav.

Man ist von den Mailändischen Ephemeriden längst gewohnt, dass sie mit einigen Aufsätzen ausgestattet sind, welche ihnen einen über das Jahr ihrer Erscheinung ausreichenden Werth geben und sie auch für solche Personen interessant machen, die nicht in dem Fall sind, von dem astronomischen Kalender Gebrauch zu machen. Von letzterm brauchen wir nichts zu sagen, als dass seine beifallswürdige Einrichtung unverändert geblieben ist: der Anhang enthält folgende Aufsätze.

Geographische Lage einiger von Mailand aus sichtbarer Berge von \textsc{Barnabas Oriani}. Dieser hochverdiente Astronom theilt uns hier ein Verzeichniss von 49 Punkten in Oberitalien und der Schweiz mit, welches ihre Breite und Länge, ihre relative Lage gegen die Cathedrale von Mailand, und ihre Höhe über der Meeresfläche enthält. Diese Bestimmungen sind grösstentheils auf die Messungen gegründet, welche die Mailänder Astronomen in den Jahren 1788---1791 auf Befehl des Oesterreichischen und in den Jahren 1803---1806 auf Befehl des Italienischen Gouvernements ausgeführt haben. Diese Messungen haben sich nicht bis zum Meere hin erstreckt, die Höhenbestimmungen sind daher nur relative gewesen, und ihrer Reduction auf die Meeresfläche ist die barometrische Höhenbestimmung von Mailand untergelegt. Bei einigen Bestimmungen des Verzeichnisses liegen Messungsoperationen aus der neuesten Zeit zum Grunde, welche hier etwas umständlicher mitgetheilt und von besonderm Interesse sind. Der Monte Viso im Piemontesischen wurde durch Messungen von Turin und Mailand aus niedergelegt, von welchem letztern Orte er 96380 Toisen (25 geographische Meilen) entfernt ist; die von beiden Orten aus gemachten Höhenmessungen vereinigen sich am besten, wenn man die terrestrische Refraction zu 0.08 der Krümmung des terrestrischen Bogens annimmt, und die Höhe über der Meeresfläche wird demzufolge von \textsc{Oriani} zu 1968 Toisen angesetzt. Die Höhenbestimmung des Monte Cimone in den Apenninen, unweit Lucca, zu 1112 Toisen, ist besonders merkwürdig, da sie sich auf unmittelbare Messung der Depression des Meereshorizonts gründet, und zwar sowohl des Mittelländischen als des Adriatischen Meeres: diese Messungen sind von \textsc{Brioschi} im Jahr 1817 ausgeführt. Die Höhe des Monte Rosa wird zu 2385, die des Finsterarhorn zu 2203, die des Simplon zu 1805 Toisen angesetzt. --- Nachrichten von den im Jahre 1822 ausgeführten Operationen, um die Längenunterschiede mehrerer Oerter in Italien, durch Pulversignale auf dem Monte Cimone, zu bestimmen, von \textsc{Francesco Carlini}. Durch die Triangulirungen in Frankreich und den Oestreichischen Staaten wird man, wenn sie vollendet und verknüpft sein werden, in den Besitz der Messung eines überaus grossen Bogens des mittlern Parallelkreises kommen, der vom Atlantischen Meere bis Orsova gegen 24 Grad betragen wird. Wie wichtig diese Messungen durch Verbindung mit zweckmässigen astronomischen Operationen für die vollkommenere Kenntniss der Gestalt der Erde werden können, fällt in die Augen. Die Beschaffenheit der Landstriche selbst, durch welche dieser Bogen geht, ist den Operationen, durch welche der Längenunterschied der Endpunkte bestimmt werden muss, besonders günstig, da auf dieser Strecke so viele hohe Berge liegen, die eine ungeheuer weite Aussicht beherrschen, so dass man mit einer verhältnissmässig sehr kleinen Anzahl von Zwischenpunkten wird ausreichen können. Ein erster Versuch dieser Art wurde schon im Sept. 1821 gemacht, indem der Längenunterschied zwischen der Mailänder Sternwarte und dem Hospiz auf dem Mont Cenis durch Pulversignale auf der 1792 Toisen hohen und 86000 Toisen von Mailand entfernten Rocca Melone bestimmt wurde. Man wünschte, aufgemuntert durch den glücklichen Erfolg dieses Versuchs, zu einer \SourceLeaf{635} umfassendern Verbindung fortzuschreiten. Die Französischen Geographen brachten dazu einen kühnen 'Plan in Vorschlag, nach welchem man vermittelst dreier Zwischenpunkte, nemlich des oben erwähnten Monte Viso, des Monte Cero bei Padua und des Monte Maggiore im Friaul in Einer Nacht die Verbindung zwischen der Ostküste des Adriatischen Meeres und des Mont d'Or bei Clermont mitten in Frankreich bewirken zu können meinte. Man stand jedoch wieder davon ab, weil man die Schwierigkeiten für zu gross hielt. Es ist noch ungewiss, ob der Monte Viso überhaupt zu ersteigen ist, noch mehr, ob man auf seiner höchsten steilen Spitze während der Nacht einen Aufenthalt machen kann. Und gesetzt auch, dass diese Schwierigkeiten sich überwinden liessen, fürchtet man, dass das Licht von Pulverblitzen bei der ungeheuern Entfernung vom Monte Cero (50 geogr. Meilen) selbst den stärksten Fernröhren unsichtbar bleiben würde (Nach diesen Aeusserungen scheint dieser riesenhafte Plan noch nicht unbedingt aufgegeben zu sein; allein Ref. findet aus den Angaben für die Höhen dieser beiden Punkte und für ihre Entfernung, dass sie gar nicht einer über den physischen Horizont des andern erhoben sein können; ohne diesen Umstand würde sich den beiden letzten Schwierigkeiten durch die Anwendung grosser Heliotrope begegnen lassen). Man entschloss sich daher einstweilen zu einer beschränktern Operation, indem man auf dem Monte Cimone im Anfang Mai 1822 mehrere Nächte hindurch Pulversignale geben liess, die auf dem Monte Cero, in Mailand und auf verschiedenen andern italienischen Sternwarten beobachtet werden sollten. Allein das ungünstige Wetter vereitelte den Erfolg dieser Operationen in der Hauptsache; weder in Mailand noch auf dem Monte Cero wurden die Pulversignale gesehen. In Parma, Modena, Bologna und Florenz wurden sie indessen beobachtet; allein die zum Theil beträchtlichen Unterschiede der Resultate von denjenigen Längendifferenzen, welche die geodätischen Operationen gegeben hatten, scheinen zu beweisen, dass die Zeitbestimmung nicht an allen diesen Orten die nöthige Genauigkeit hatte. Ganz besonders merkwürdig ist noch, dass Hr. \textsc{Carlini} aus seinen Beobachtungen mit einem 18zolligen \textsc{Reichenbach}'schen Repetitionskreise die Polhöhe von Parma um $22''6$ grösser gefunden hat, als sie sich aus der geodätischen Verbindung mit Mailand ergeben hat. Da sich ähnliche Anomalien bei mehrern andern Oertern Oberitaliens schon früher gezeigt haben, so ist es schwer, deren Realität in Zweifel zu ziehen, und es ist sehr zu wünschen, dass alle Hauptdreiecke der Oestreichischen Triangulirungen bald vollständig bekannt gemacht werden mögen. --- Die übrigen Artikel des Anhangs enthalten noch: die von \textsc{Angelo Cesaris} in den Jahren 1817 und 1818 beobachteten Oppositionen des Uranus; die von demselben Astronomen beobachteten Oppositionen des Jupiter und Saturn im Jahr 1821; beobachtete Sternbedeckungen und Jupiters-Trabanten-Verfinsterungen von \textsc{Hallaschka} in Prag; Beobachtungen des ersten Cometen von 1822 von demselben; endlich die meteorologischen Beobachtungen in Mailand vom Jahre 1820 von A. \textsc{Cesaris}.
